
\section{Non-stationary TGP equations of motion}

\subsection{Starting point: the time-dependent field equation}

The TGP field $\Phi(x,t)$ obeys the Euler--Lagrange equation derived from
the action $S = \int dt\, E[\Phi]$ extended with a kinetic term.
Including the standard minimal kinetic term for the scalar field gives the
field equation
\begin{equation}
\Box\,\Phi + V'(\Phi;\Phi_0)
= \sum_{i=1}^N C_i\,\delta^{(3)}\!\bigl(x - x_i(t)\bigr),
\label{eq:field_eq_full}
\end{equation}
where $\Box = \partial_t^2 - \nabla^2$ is the d'Alembertian and
the source $C_i\,\delta^{(3)}$ models each body as a point source of
field strength $C_i$.  The self-interaction potential is
\begin{equation}
V(\Phi;\Phi_0)
= \frac{\beta}{3\Phi_0}\,\Phi^3 - \frac{\gamma}{4\Phi_0^2}\,\Phi^4.
\end{equation}

\paragraph{Vacuum condition.}
The field sits in the non-trivial vacuum $\Phi_0>0$ with $\beta=\gamma$
(exact TGP vacuum condition), giving a screening mass
\begin{equation}
m_{\rm sp}^2 = V''(\Phi_0;\Phi_0)\big|_{\beta=\gamma} = \beta.
\end{equation}

\subsection{Linearisation around the vacuum}

Write $\Phi(x,t) = \Phi_0\bigl(1 + \delta(x,t)\bigr)$ with $\delta\ll 1$.
Expanding $V'$ to linear order in $\delta$:
\begin{equation}
V'(\Phi) \approx V'(\Phi_0) + V''(\Phi_0)\,\Phi_0\,\delta
= m_{\rm sp}^2\,\Phi_0\,\delta,
\end{equation}
where we used $V'(\Phi_0)=0$ (vacuum equation) and $V''(\Phi_0)=m_{\rm sp}^2$.
Hence
\begin{equation}
\boxed{
\bigl(\Box - m_{\rm sp}^2\bigr)\delta(x,t)
= \frac{1}{\Phi_0}\sum_{i=1}^N C_i\,\delta^{(3)}\!\bigl(x - x_i(t)\bigr).
}
\label{eq:wave_delta}
\end{equation}
This is the \emph{Klein--Gordon equation with time-dependent point sources}.

\subsection{Retarded solution: the radiation field}

The causal solution of \eqref{eq:wave_delta} is given by convolution with
the retarded Green function $G_{\rm ret}$:
\begin{equation}
\bigl(\Box - m_{\rm sp}^2\bigr)G_{\rm ret}(x,t;x',t')
= \delta^{(3)}(x-x')\,\delta(t-t'),
\end{equation}
whose explicit form in 3+1 dimensions is (Yukawa propagator):
\begin{equation}
G_{\rm ret}(x,t;x',t')
=
-\frac{1}{4\pi |x-x'|}
\left[
  \delta\!\bigl(t-t'-|x-x'|\bigr)
  - m_{\rm sp}\,\frac{J_1\!\bigl(m_{\rm sp}\sqrt{(t-t')^2-|x-x'|^2}\bigr)}
                     {\sqrt{(t-t')^2-|x-x'|^2}}\,
    \theta\!\bigl(t-t'-|x-x'|\bigr)
\right].
\end{equation}

Therefore the exact linearised field is
\begin{equation}
\delta(x,t)
=
\frac{1}{\Phi_0}
\sum_{i=1}^N C_i
\int_{-\infty}^{t} dt'\;
G_{\rm ret}\!\bigl(x,t;\, x_i(t'),\, t'\bigr).
\end{equation}

\paragraph{Physical interpretation.}
The key difference from the quasi-static case is that:
\begin{itemize}
  \item each body emits a \emph{wave} that propagates at speed $c=1$
    (in natural units);
  \item the wave decays exponentially on length scale $m_{\rm sp}^{-1}$
    \emph{in the rest frame};
  \item the field at body $j$ at time $t$ depends on the \emph{retarded
    position} $x_i(t - |x_j - x_i|)$, not the current position.
\end{itemize}

\subsection{Non-stationary pairwise interaction}

For two slowly moving bodies (velocities $v \ll 1$), the retarded position
can be expanded:
\begin{equation}
x_i\!\bigl(t - d_{ij}/c\bigr)
= x_i(t) - \dot x_i(t)\,d_{ij}/c + \tfrac{1}{2}\ddot x_i(t)\,(d_{ij}/c)^2 + \cdots
\end{equation}
This gives corrections to the static Yukawa potential of order $v/c$.

The leading-order retardation correction to the pairwise energy is (in the
$m_{\rm sp} d\ll 1$ limit):
\begin{equation}
\delta V_2^{(v/c)}
= \frac{4\pi C_i C_j}{d_{ij}}\cdot v_i \cdot v_j / c^2 + O(v^2/c^2),
\end{equation}
analogous to the Darwin term in electrodynamics.

\subsection{Radiation: power loss}

For an accelerating source (e.g.\ body $i$ with acceleration $\ddot x_i$),
the scalar field radiates energy.  In the far field ($r\to\infty$, $m r\ll 1$),
the radiated power is (Larmor formula for massive scalar field):
\begin{equation}
P_{\rm rad,i}
= \frac{C_i^2}{6\pi}\,|\ddot x_i|^2
  \cdot
  \begin{cases}
    1 & m_{\rm sp}=0\text{ (massless)},\\[4pt]
    \displaystyle
    1 - \frac{3m_{\rm sp}^2}{4\omega^2}+\cdots & \omega \gg m_{\rm sp},
  \end{cases}
\end{equation}
where $\omega$ is the orbital frequency.  For TGP screening masses
$m_{\rm sp} = \sqrt{\beta}$ and typical orbital frequencies,
this quantifies the rate of energy loss to scalar radiation.

\subsection{Non-stationary N-body equations of motion}

Including retardation and radiation reaction, the equation of motion for
body $i$ becomes:
\begin{equation}
\boxed{
C_i\,\ddot x_i(t)
=
\underbrace{
F_i^{(2,\rm quasi)}\!\bigl(x_i(t),\{x_j(t)\}\bigr)
}_{\text{quasi-static pairwise}}
+
\underbrace{
F_i^{(3,\rm quasi)}\!\bigl(\{x_j(t)\}\bigr)
}_{\text{quasi-static 3-body}}
+
\underbrace{
\delta F_i^{(v/c)}\!\bigl(\{x_j(t),\dot x_j(t)\}\bigr)
}_{\text{velocity corrections}}
+
\underbrace{
\delta F_i^{(\rm RR)}(t)
}_{\text{radiation reaction}}
+
\cdots
}
\label{eq:eom_nonstationary}
\end{equation}

The terms are:
\begin{enumerate}
\item \emph{Quasi-static pairwise force} (exact, closed form):
\begin{equation}
F_i^{(2,\rm quasi)}
= -\nabla_{x_i}\sum_{j\ne i} V_2(d_{ij}),\quad
V_2(d) = -\frac{4\pi C_i C_j}{d} + \frac{8\pi\beta C_i C_j}{d^2}
         - \frac{12\pi\gamma C_i C_j(C_i+C_j)}{d^3}.
\end{equation}

\item \emph{Quasi-static 3-body force} (exact via Feynman integral):
\begin{equation}
F_i^{(3,\rm quasi)}
= 6\gamma\sum_{j<k\atop j\ne i,k\ne i} C_i C_j C_k
\left[
\frac{\partial I_Y}{\partial d_{ij}}\,\hat e_{ij}
+ \frac{\partial I_Y}{\partial d_{ik}}\,\hat e_{ik}
\right],
\end{equation}
where $\partial I_Y/\partial d_{ij}$ is computed exactly via
\begin{equation}
\frac{\partial I_Y}{\partial d_{12}}
=
-2m\,d_{12}
\int_{\Delta_2}
\frac{\alpha_2}{\Delta^2 \sqrt{Q}}\,K_1\!\Bigl(m\sqrt{Q/\Delta}\Bigr)\,d\alpha,
\label{eq:exact_dI_dd}
\end{equation}
with $Q = \alpha_2 d_{12}^2 + \alpha_1 d_{13}^2 + \alpha_3 d_{23}^2$
and $\Delta = \alpha_1\alpha_2+\alpha_1\alpha_3+\alpha_2\alpha_3$
(implemented in \texttt{three\_body\_force\_exact.py}).

\item \emph{Velocity corrections} ($v/c$ order): retardation shifts, Darwin term,
  Römer delay in the effective potential.  Suppressed by $(v/c)^2\sim (v_{\rm orb}/1)^2$.

\item \emph{Radiation reaction} (Abraham--Lorentz--Dirac term for a scalar):
\begin{equation}
\delta F_i^{(\rm RR)}
= \frac{C_i^2}{6\pi}\,\dddot x_i(t).
\end{equation}
  Suppressed by $C_i^2$ (source strength squared) and relevant only when
  the orbital period is comparable to the radiation timescale.
\end{enumerate}

\subsection{Hierarchy of approximations}

\begin{center}
\begin{tabular}{lll}
\hline
Level & Approximation kept & Error scale\\
\hline
0 & Full field equation \eqref{eq:field_eq_full} (nonlinear, retarded) & exact\\
1 & Linearised + retarded \eqref{eq:wave_delta} & $O(\delta^2)$\\
2 & Quasi-static: drop $\partial_t^2\delta$ & $O(v^2/c^2)$\\
3 & Yukawa superposition: $\delta=\sum \delta_i$ & $O(\delta^2)$\\
4 & Truncate at 3-body: drop $V_4,V_5,\ldots$ & $O(C^4)$\\
5 & Exact 3-body force via \eqref{eq:exact_dI_dd} & $0$ (within level 4)\\
\hline
\end{tabular}
\end{center}

\paragraph{What \texttt{dynamics\_v2.py} implements.}
Currently levels 0--4 are applied and, for the 3-body term, an
additional saddle-point approximation $I_Y \approx (2\pi)^{3/2}m^{-2}
e^{-ms}/s$ is used (roughly 5--30\% error in the 3-body force).

\paragraph{What \texttt{three\_body\_force\_exact.py} provides.}
The exact 3-body force at level 5, replacing the saddle-point
approximation with the exact Feynman 2D integral, at the cost of
$O(n_{\rm quad}^2)$ function evaluations per triplet per time step.

\subsection{Retardation time scale}

The quasi-static approximation (level 2) fails when the retardation
delay $\tau = d_{ij}/c$ is comparable to the orbital period $T_{\rm orb}$.
The condition for validity is
\begin{equation}
\frac{d_{ij}}{c\,T_{\rm orb}} = \frac{v_{\rm orb}}{c} \ll 1.
\end{equation}
For TGP bodies moving at non-relativistic speeds ($v_{\rm orb}/c \sim
10^{-3}$--$10^{-6}$ in astrophysical contexts), the quasi-static limit is
excellent.  Post-Newtonian corrections enter at relative order $v^2/c^2
\sim 10^{-6}$--$10^{-12}$ and are negligible for most TGP applications.

\subsection{Non-stationary 3-body effects: summary}

The genuine novelty of the non-stationary TGP equations (compared to
Newtonian gravity or standard Yukawa) lies in two effects:

\begin{enumerate}
\item \emph{Non-trivial 3-body interaction} (already present in the
  quasi-static limit): the term $V_3 \propto \gamma\,C_1 C_2 C_3\,I_Y$
  has no Newtonian analogue.  It is attractive, breaks the equivalence
  with gravity ($C_i\ne$ mass), and is dominant at intermediate separations
  $m_{\rm sp}d\sim 1$.

\item \emph{Modified radiation}: the massive scalar ($m_{\rm sp}>0$) radiates
  less efficiently than a massless field.  The radiation spectrum has a
  low-frequency cutoff at $\omega = m_{\rm sp}$, unlike electromagnetic or
  gravitational radiation.  This means short-period orbits are more
  stable in TGP than in GR.
\end{enumerate}

The leading observable signatures of non-stationary effects in TGP
would therefore be:
\begin{itemize}
  \item orbital decay rate suppressed below the gravitational-wave analogue
    by a factor of $(m_{\rm sp}\,d_{\rm orb})^2$ relative to GR,
  \item three-body secular effects (Kozai--Lidov-type resonances modified
    by the $V_3$ coupling) measurable at the $C^3/(C^2)=C$ level compared
    to pairwise dynamics.
\end{itemize}
